\documentclass[12pt,letterpaper]{article}

\usepackage[utf8]{inputenc}
\usepackage[T1]{fontenc}
\usepackage[spanish,es-nodecimaldot]{babel}
\usepackage{lmodern}
\usepackage{microtype}
\usepackage{amsmath,amssymb}
\usepackage{booktabs,tabularx,array}
\usepackage{geometry}
\usepackage{fancyhdr}
\usepackage{hyperref}

\geometry{left=2.5cm,right=2.5cm,top=2.6cm,bottom=2.6cm,headheight=16pt}
\setlength{\parindent}{0pt}
\setlength{\parskip}{0.55em}
\renewcommand{\arraystretch}{1.18}
\newcolumntype{Y}{>{\raggedright\arraybackslash}X}
\newcommand{\code}[1]{\texttt{#1}}
\newcommand{\q}{Q-learning}

\pagestyle{fancy}
\fancyhf{}
\fancyhead[L]{\small Finanzas Descentralizadas IV}
\fancyhead[R]{\small Parcial 2 -- Q-learning}
\fancyfoot[C]{\thepage}
\renewcommand{\headrulewidth}{0.4pt}

\title{\textbf{Del modelo de Bellman a Q-learning}\\[0.4em]
\large Agente de decisión POL/USDC guiado por evidencia de wallets on-chain}
\author{Juan Camilo Puentes y Jesús Figueroa\\
Finanzas Descentralizadas IV\\
Profesor: Julián Sánchez}
\date{Parcial 2 -- Septiembre de 2026}

\begin{document}
\maketitle

\begin{abstract}
Este documento continúa el proyecto del Parcial 1 sobre decisiones de cartera
entre USDC y POL. El agente observa el régimen reciente de POL, una señal
causal construida con wallets de desempeño consistente y la posición actual.
La formulación original estimaba un MDP y resolvía Bellman hacia atrás. En este
avance se conserva el mismo estado, las mismas acciones y la recompensa neta de
gas, pero se aprende una tabla temporal $Q_t(s,a)$ mediante experiencias
históricas y exploración $\varepsilon$-greedy. Se presenta una evaluación
cronológica separada: 16 horas para entrenar y 8 horas posteriores para probar.
El resultado ilustra el mecanismo y sus límites; no constituye una afirmación de
rentabilidad futura.
\end{abstract}

\section{Problema y decisión del agente}

El problema es administrar, cada hora, una cartera simplificada con capital en
USDC o POL. El \textbf{agente} es la política de aprendizaje por refuerzo. El
\textbf{entorno} es el historial del pool Uniswap V3 POL/USDC en Polygon,
junto con la actividad on-chain de las wallets. La \textbf{decisión secuencial}
consiste en elegir una acción válida para la posición disponible y recibir el
resultado observado durante la hora siguiente.

El objetivo financiero es maximizar la riqueza acumulada neta de gas. Las
wallets no se copian como una regla fija: sus perfiles agregados sólo aportan
contexto al estado y un ajuste prudente de la recompensa cuando su señal coincide
con una decisión financieramente ventajosa.

\begin{center}
\begin{tabularx}{0.92\linewidth}{p{3cm}Y}
\toprule
\textbf{Componente} & \textbf{Definición en el proyecto}\\
\midrule
Agente & Política tabular que elige acciones sobre una cartera USDC/POL.\\
Entorno & Precios horarios del pool POL/USDC y señales causales de wallets on-chain.\\
Frecuencia & Una decisión por hora.\\
Episodio & Secuencia ordenada de cortes horarios; el snapshot usado contiene 24 observaciones.\\
Objetivo & Aprender el valor esperado de cada acción válida, descontando gas y considerando el efecto futuro de la cartera.\\
\bottomrule
\end{tabularx}
\end{center}

\section{Del MDP de Bellman a la formulación de Q-learning}

En Bellman se estimaron recompensas medias y probabilidades de transición
$P(s'\mid s,a)$ para resolver la política de horizonte finito hacia atrás. En
\q\ no se consulta la matriz $P$ durante la actualización. En su lugar, el
agente recorre una experiencia observada, compara lo que esperaba con lo que
ocurrió y corrige una sola entrada de su tabla.

La formulación concreta es
\[
S_t=(R_t,W_t,X_t),\qquad A_t\in\mathcal A(S_t),\qquad
R_{t+1}=r(S_t,A_t,S_{t+1}).
\]

\begin{center}
\begin{tabularx}{0.95\linewidth}{p{2.6cm}Y Y}
\toprule
\textbf{Variable} & \textbf{Qué representa} & \textbf{Qué comportamiento captura}\\
\midrule
$R_t$ & Régimen: \code{DOWN}, \code{FLAT} o \code{UP}. & Promedio de retornos pasados de 1m, 5m, 15m y 1h; no usa precio futuro.\\
$W_t$ & Señal: \code{BUY}, \code{SELL}, \code{HOLD} o \code{NEUTRAL}. & Consenso de perfiles de 1h que ya maduraron al corte $t$.\\
$X_t$ & Posición: 0 para USDC y 1 para POL. & Exposición al precio y acciones legalmente disponibles.\\
$A_t$ & Acción del agente. & Entrar, salir o conservar la posición.\\
$R_{t+1}$ & Recompensa observada. & Retorno de POL, gas y apoyo de wallets.\\
\bottomrule
\end{tabularx}
\end{center}

El espacio tiene $3\times4\times2=24$ estados interpretables. La
representación es una aproximación de Markov: resume información útil sin
guardar todo el historial. Su principal limitación es que no incluye liquidez,
slippage, volumen del pool, cambios de volatilidad de largo plazo ni balances
externos de las wallets. Una ampliación razonable es añadir una categoría de
volatilidad o liquidez sin abandonar el modelo tabular.

\subsection{Acciones admisibles}

El agente es long-only y no usa deuda ni apalancamiento. Por ello una acción
debe ser compatible con la posición actual.

\begin{center}
\begin{tabular}{lll}
\toprule
\textbf{Posición actual} & \textbf{Acciones admisibles} & \textbf{Posición siguiente}\\
\midrule
USDC & \code{BUY\_POL}, \code{HOLD} & POL o USDC\\
POL & \code{SELL\_POL}, \code{HOLD} & USDC o POL\\
\bottomrule
\end{tabular}
\end{center}

\section{Recompensa y experiencia de aprendizaje}

Sea $g_t=(p_{t+1}/p_t)-1$ el retorno observado de POL en una hora y $c_t$ el
costo de gas relativo. La recompensa base $b_t$ es
\[
b_t(s,a)=
\begin{cases}
g_t-c_t, & X_t=0,\ a=\code{BUY\_POL},\\
0, & X_t=0,\ a=\code{HOLD},\\
-c_t, & X_t=1,\ a=\code{SELL\_POL},\\
g_t, & X_t=1,\ a=\code{HOLD}.
\end{cases}
\]

La señal de wallets puede reforzar la acción dominante sólo cuando coincide con
una ventaja financiera. Si $\hat a_t^W$ es la acción respaldada, $q_t^W$ la
confianza de la señal y $\bar a$ la alternativa admisible, el aprendizaje usa
\[
R_{t+1}=b_t(s,a)+0.25\,q_t^W\,
\mathbf{1}\{a=\hat a_t^W\}\,[b_t(s,a)-b_t(s,\bar a)].
\]
El primer término es siempre financiero. El segundo no convierte una mala
operación en buena: si la acción apoyada tiene peor recompensa base, su ventaja
es negativa. Para comparar desempeño con Buy and Hold se reporta por separado
la riqueza económica, calculada sólo con $b_t$.

\subsection{Una transición explicada}

Considérese una transición simulada coherente con el proyecto. El agente está
en
\[
s_t=(\code{UP},\code{BUY},\text{USDC}),\qquad a_t=\code{BUY\_POL}.
\]
POL pasa de 1.00 a 1.02 USDC, el gas relativo es $0.001$ y la confianza Buy es
$0.80$. La recompensa base es $0.02-0.001=0.019$. La ventaja frente a esperar
en USDC es $0.019$, por lo que el ajuste de wallets es
$0.25(0.80)(0.019)=0.0038$. Entonces $R_{t+1}=0.0228$ y el estado siguiente
puede ser $s_{t+1}=(\code{FLAT},\code{HOLD},\text{POL})$.

Si para las acciones admisibles de $s_{t+1}$ el máximo aprendido es $0.0300$,
el target con $\gamma=0.99$ es
\[
Y_t=0.0228+0.99(0.0300)=0.0525.
\]
Supóngase que antes de la experiencia $Q_t(s_t,\code{BUY\_POL})=0.0100$ y que
$\alpha=0.15$. La actualización queda
\[
Q_t(s_t,a_t)\leftarrow0.0100+0.15(0.0525-0.0100)=0.0164.
\]
La experiencia fue mejor de lo que la tabla esperaba, por lo que aumenta el
valor de comprar en ese contexto. Si el target fuera menor que el valor previo,
la tabla reduciría su estimación. Con datos cortos, un $\alpha$ muy grande
haría que cada hora atípica tenga demasiado peso.

\section{Exploración, explotación y entrenamiento}

La tabla se inicializa en cero con una dimensión por hora restante:
$Q_t(s,a)$ para $t=0,\ldots,15$ en el bloque de entrenamiento. Se entrenó con
2,000 episodios sobre las primeras 16 observaciones del snapshot. Cada episodio
recorre el bloque en orden temporal, parte de USDC y actualiza sólo acciones
admisibles.

La política de comportamiento es $\varepsilon$-greedy. En el episodio $e$,
$\varepsilon$ disminuye linealmente de 0.90 a 0.05:
\[
\varepsilon_e=0.90+\frac{e-1}{2000-1}(0.05-0.90).
\]
\textbf{Explorar} significa escoger al azar entre las dos acciones válidas.
\textbf{Explotar} significa escoger la acción admisible con el mayor valor Q.
El valor inicial alto evita que un empate con Q inicial cero fije toda la
trayectoria; el piso de 0.05 conserva exploración hasta el final. En evaluación
se fija $\varepsilon=0$: no se actualiza Q ni se sortean acciones.

\section{Evidencia de aprendizaje}

La evidencia no se limita a la riqueza final. La Tabla~\ref{tab:evidencia}
muestra que la exploración se reduce, la recompensa por episodio aumenta en la
trayectoria de entrenamiento y la magnitud media de Q se estabiliza. La
recompensa acumulada depende de las acciones exploradas, por lo que no tiene que
crecer de forma monótona.

\begin{table}[h]
\centering
\caption{Diagnóstico de entrenamiento de Q-learning en las 16 horas iniciales.}
\label{tab:evidencia}
\begin{tabular}{rrrr}
\toprule
\textbf{Episodio} & $\boldsymbol{\varepsilon}$ & \textbf{Recompensa total} & \textbf{Media $|Q|$}\\
\midrule
1 & 0.9000 & 0.1329 & 0.000030\\
500 & 0.6878 & 0.0619 & 0.004037\\
1000 & 0.4752 & 0.1299 & 0.004039\\
1500 & 0.2626 & 0.2612 & 0.004039\\
2000 & 0.0500 & 0.3151 & 0.004039\\
\bottomrule
\end{tabular}
\end{table}

El promedio de recompensa de los primeros 200 episodios fue 0.1062 y el de los
últimos 200 fue 0.2916. Esto muestra que la tabla cambia y la política explota
más las acciones con mayor valor. No demuestra una ganancia estable fuera de
muestra, porque el bloque de entrenamiento es corto y repite una trayectoria.

\section{Política aprendida e interpretación}

La política greedy se obtiene mediante
\[
\pi_t(s)\in\arg\max_{a\in\mathcal A(s)}Q_t(s,a).
\]
La Tabla~\ref{tab:politica} presenta estados que aparecieron durante el bloque
de entrenamiento. Los valores cambian con $t$ porque quedan distintas horas por
decidir; el mismo estado puede tener otra acción en una etapa posterior.

\begin{table}[h]
\centering
\caption{Estados representativos de la política Q aprendida.}
\label{tab:politica}
\begin{tabularx}{\linewidth}{p{1.2cm}YrrY}
\toprule
$t$ & \textbf{Estado} & \textbf{Q de acción 1} & \textbf{Q de acción 2} & \textbf{Política e interpretación}\\
\midrule
0 & $(\code{FLAT},\code{BUY},\text{USDC})$ & Buy: 0.3029 & Hold: 0.2543 & \code{BUY\_POL}. Al inicio, la combinación observada favorece entrar a POL.\\
1 & $(\code{UP},\code{HOLD},\text{POL})$ & Sell: 0.1942 & Hold: 0.2596 & \code{HOLD}. Conserva la exposición cuando el valor futuro de mantener supera al de salir.\\
3 & $(\code{UP},\code{HOLD},\text{POL})$ & Sell: 0.1457 & Hold: 0.1150 & \code{SELL\_POL}. Con menos tiempo restante, el valor favorece realizar la salida.\\
8 & $(\code{DOWN},\code{BUY},\text{USDC})$ & Buy: 0.0213 & Hold: 0.0480 & \code{HOLD}. Aunque la señal es Buy, el retorno y horizonte aprendidos favorecen esperar.\\
\bottomrule
\end{tabularx}
\end{table}

El cuarto caso muestra por qué la señal de wallets no es una orden automática.
Los conteos de visitas del estado inicial fueron 1,554 para Buy y 446 para Hold,
de un total de 2,000 episodios. Esos conteos ayudan a interpretar la confianza
de la tabla, pero no resuelven la falta de cobertura en estados no observados.

\section{Evaluación cronológica y comparación}

El Parcial 2 exige evaluar con información distinta de la usada para aprender
Q. Se hizo una división cronológica sin mezclar horas:

\begin{center}
\begin{tabular}{lll}
\toprule
\textbf{Bloque} & \textbf{Cortes horarios UTC} & \textbf{Uso}\\
\midrule
Entrenamiento & 22-Ago-2026 02:00 a 17:00 & Actualizar Q.\\
Evaluación & 22-Ago-2026 18:00 a 23-Ago-2026 01:00 & Política congelada, sin actualizar Q.\\
\bottomrule
\end{tabular}
\end{center}

La métrica principal es la riqueza económica con capital inicial 100, compuesta
con recompensa base y por tanto neta de gas. Se compara con Bellman entrenado en
el mismo bloque y con Buy and Hold, que compra POL una vez al inicio de la
evaluación y mantiene hasta el final.

\begin{table}[h]
\centering
\caption{Resultado fuera del bloque de entrenamiento: riqueza económica neta de gas.}
\begin{tabular}{lrrl}
\toprule
\textbf{Estrategia} & \textbf{Riqueza final} & \textbf{Retorno} & \textbf{Acciones durante evaluación}\\
\midrule
Q-learning & 97.9904 & -2.01\% & 4 compras y 4 ventas\\
Bellman & 97.5524 & -2.45\% & 1 compra, 1 venta y 6 hold\\
Buy and Hold & 96.2086 & -3.79\% & compra inicial y mantenimiento\\
\bottomrule
\end{tabular}
\end{table}

Q-learning obtuvo la menor pérdida en esta única ventana de evaluación, pero no
supera una prueba estadística ni permite concluir superioridad. Varios estados
del bloque de prueba no fueron visitados en entrenamiento, por lo que sus
valores Q permanecen cercanos a cero y los empates pueden determinar la acción.
El resultado negativo es informativo: la implementación se puede evaluar y la
generalización todavía es limitada.

\section{Interpretación, limitaciones y siguiente avance}

El agente aprendió valores distintos para acciones admisibles en estados que sí
recorrió durante el entrenamiento, y su decisión puede variar con las horas que
quedan. Parece funcionar mejor en combinaciones repetidas, como
$(\code{UP},\code{BUY},\text{POL})$, donde ambas alternativas recibieron muchas
actualizaciones. Presenta dificultades en estados escasos o ausentes y en una
ventana de sólo 24 horas con señales de wallets concentradas en pocas categorías.

La principal limitación es el tamaño y diversidad de la historia. Una tabla
discreta necesita visitas para estimar cada par estado-acción; con pocos cortes,
la división temporal deja muy poca cobertura en prueba. También faltan slippage,
liquidez, comisiones variables e información de mercado más amplia.

El siguiente avance concreto es acumular snapshots de múltiples fechas,
reservar bloques cronológicos para entrenamiento, validación y prueba, y medir
la cobertura de cada estado antes de comparar estrategias. Con esa historia se
podrán ajustar $\alpha$, $\gamma$ y el decaimiento de $\varepsilon$ sólo en
validación y conservar el test final intacto.

\section{Conclusión}

La transición desde Bellman a Q-learning conserva la semántica del proyecto:
el estado integra régimen, wallets y posición; las acciones respetan una cartera
long-only; y la recompensa sigue siendo financiera y neta de gas. La diferencia
es metodológica: Bellman usa un modelo estimado del entorno, mientras Q-learning
ajusta sus valores con experiencias.

La evaluación cronológica evidencia que el pipeline ya separa aprendizaje y
prueba, aunque el volumen de datos no basta para una conclusión financiera. El
resultado cumple el objetivo del Parcial 2: mostrar qué aprende el agente, cómo
explora, cómo se actualiza Q y cómo interpretar críticamente su desempeño.

\begin{thebibliography}{9}
\bibitem{parcial1}
Puentes, J. C. y Figueroa, J. (2026). \emph{Perfiles de wallets on-chain y programación dinámica para decisiones de operación sobre POL}. Parcial 1, Finanzas Descentralizadas IV.

\bibitem{parcial2}
Sánchez, J. (2026). \emph{Parcial 2 -- Sustentación del proyecto: del modelo de Bellman a Q-learning}. Finanzas Descentralizadas IV.

\bibitem{sutton}
Sutton, R. S. y Barto, A. G. (2018). \emph{Reinforcement Learning: An Introduction} (2.a ed.). MIT Press.
\end{thebibliography}

\end{document}
